\documentclass[utf8,a4paper,12pt,headsepline,oneside,numbers=noenddot]{scrbook}
\usepackage[utf8]{inputenc}
\usepackage{url}
\usepackage{geometry}
\usepackage[font=small,labelfont=bf]{caption}
%\usepackage[ansinew]{inputenc}
\usepackage{calc}
\usepackage{graphicx}
\usepackage[table,xcdraw]{xcolor}
\usepackage{tabularx}
\usepackage{nameref}
\usepackage{afterpage}
%\usepackage{floatflt}
\usepackage{listings}
\usepackage{wrapfig}
\usepackage{amsmath, amsthm, amssymb}
\usepackage{mathtools}
%\usepackage{longtable} %Support for large tables
%\usepackage{rotating} %Support for rotation of large tables
\usepackage{makeidx}
%\usepackage[style=altlist, number=none, toc]{glossary} %Glossary support
\usepackage{algorithm}
\usepackage{listings}%[procnames]
\usepackage[nottoc]{tocbibind}
\usepackage[ngerman]{babel}
\usepackage{scrpage2}
\usepackage{textcomp}
\usepackage{setspace} % Achtung: unterhalb von \begin{document} wird der Zeilenabstand variabler gesetzt (1.5 absoluter MAX - Wert!!!)
\usepackage{breakcites} % Bricht zu lange Zitierungen im Text um (vorallem benötigt by apalike,chicago, ....)
\usepackage{titlesec}
\usepackage{geometry} 
\usepackage{array, makecell}

% decrease vspace on chapter page



\titleformat{\chapter}{\normalfont \Large \bfseries }{ \thechapter}{20pt}{\LARGE\bfseries}

\titlespacing*{\chapter}{0pt}{-10pt}{20pt}

%Set page layout
\geometry{verbose,a4paper,left=3cm,top=2.5cm+\headheight+\headsep,right=2.5cm,bottom=2.5cm+\footskip,pdftex}

%Additional latex definitions

% Definition of a tight list (for title page)
\newenvironment{tightlist}{
\begin{list}{\textbullet}{
\setlength{\topsep}{-1ex}
\setlength{\itemsep}{-1ex}
\setlength{\leftmargin}{4ex}
}
}{
\end{list}
\vspace{1ex}
}

%Default font
\renewcommand{\familydefault}{\sfdefault}

%Create glossary, index

%\makeglossary
\makeindex

%-Eingabe der Metadaten des Titelblattes--------------------------

%-Daten des Autors / Authors Data---------------------------------

\newcommand{\dcauthorpre}{Herr} %ggf Titel
\newcommand{\dcauthorsurname}{Schack} 
\newcommand{\dcauthorname}{Sebastian} 
\newcommand{\dcauthoradd}{29.10.1991} %geboren am 
\newcommand{\dcauthorno}{6055} %Matrikelnummer

%-Titel und Untertitel / Title and subtitle-----------------------

\newcommand{\dctitle}{Analyse und Neukonzeption eines Kreditorenworkflows auf Basis von Vorgehensmodellen für die Gestaltung von Unternehmensarchitekturen am Beispiel der GLC Glücksburg Consulting AG} 
\newcommand{\dcsubtitle}{~}  
% Falls dcsubtitle NICHT verwendet werden soll, {\dcsubtitle}{~} eingeben.

%-Eingabe der Gutachternamen / Names of the approvals-------------

\newcommand{\dcapprovala}{Dr.-Ing. Jan Himmelspach} 
\newcommand{\dcapprovalb}{Dr. Hinrich Schröder} 
\newcommand{\dcapprovalc}{} 

%-Eingabe der Betreuernamen / Name of the advisors ---------------

\newcommand{\dcada}{Peter Weigel}

%-Information zur Universitaet------------------------------------

\newcommand{\dcdegree}{Bachlor of Science\\(BSc)} 
\newcommand{\dcsubject}{Informatik} 
\newcommand{\dcfaculty}{Fachbereich Informatik}
\newcommand{\dcuniversity}{NORDAKADEMIE}
\newcommand{\dcdean}{}
\newcommand{\dcpresident}{Prof. Dr. Plate} %???

%-Pruefungsdaten: eingereicht und mdl. Pruefung-------------------
%-data of submission and oral exam--------------------------------

\newcommand{\dcdatesubmitted}{31. Januar 2011} %auch wenn nicht auf dem Titelblatt, bitte erfüllen!
\newcommand{\dcdateexam}{n.n.} 

%-CR Klassifikation / cr classification --------------------------
%siehe: http://www.acm.org/class/1998/
\newcommand{\dccrclassification}{I.6.8[Simulation and Modelling]:Types of Simulation---Discrete event, Distributed, Parallel;}

%-deutsche Schlagwoerter / german keywords------------------------

\newcommand{\dckeydea}{Schlagwort 1}
\newcommand{\dckeydeb}{Schlagwort 2}
\newcommand{\dckeydec}{Schlagwort 3}
\newcommand{\dckeyded}{Schlagwort 4}

% Folgende Zeile bitte nicht aendern!
\newcommand{\dckeywordsde}{\vfill \raggedright {\textbf{Schlagw\"orter:}}\\ \dckeydea, \dckeydeb, \dckeydec, \dckeyded \\}

%-englische Schlagwoerter / english keywords----------------------

\newcommand{\dckeyena}{keyword 1}
\newcommand{\dckeyenb}{keyword 2}
\newcommand{\dckeyenc}{keyword 3}
\newcommand{\dckeyend}{keyword 4}

% Folgende Zeile bitte nicht aendern!
\newcommand{\dckeywordsen}{\vfill \raggedright {\textbf{Keywords:}}\\ \dckeyena, \dckeyenb, \dckeyenc, \dckeyend \\}
%\newcommand{\dccrandkeywordsen}{\vfill \raggedright { \small {\textbf{CR--Classification (1998)::}}\\ \dccrclassification \\}\vspace{4mm}{\textbf{Keywords:}}\\ \dckeyena, \dckeyenb, \dckeyenc, \dckeyend \\}
%\raggedright removed from line below (after \vfill)
\newcommand{\dccrandkeywordsen}{\vfill { \small \noindent {\textbf{CR--Classification (1998)::}}\\ \dccrclassification \\}\vspace{4mm}{\textbf{Keywords:}}\\ \dckeyena, \dckeyenb, \dckeyenc, \dckeyend \\}

\newcommand{\dcpdfsubject}{Bachelorarbeit}
							% Bitte ALLE Angaben erfüllen!
\usepackage{ifpdf}

% % Um dieses Document sowohl fuer latex als auch fuer pdflatex
% % kompatibel zu machen, definieren wir uns die Variable ifpdf
%%% check whether we are running pdflatex
%\newif\ifpdf
%\ifx\pdfoutput\undefined
%\pdffalse 		% we are not running pdflatex
%\else
%\pdfoutput=1 		% we are running pdflatex
%%das muessen wir nicht unbedingt uebernehmen
%\pdfcompresslevel=9    % compression level for text and image;
%\pdftrue 
%\fi

\ifpdf
%%das kann man benutzen, wenn man andere Formate benutzen will
%\DeclareGraphicsExtensions{{.pdf}}   %Endung der Grafiken, wenn nicht pdf
% die folgenden Angaben sind im PDF unter Datei | Dokumenteigenschaften 
% in Acrobat / Acrobat Reader sichtbar
% Aendern Sie bitte die Daten, wo noetig!
\usepackage[%
	pdftitle={\dctitle},
	pdfauthor={\dcauthorsurname\ \dcauthorname},
	pdfsubject={\dcpdfsubject}, % optional
	pdfkeywords={\dckeydea, \dckeydeb, \dckeydec, \dckeyded},
	pdfpagemode=UseOutlines,
  colorlinks=true,					% bitte nicht ändern!
	linkcolor=black,					% bitte nicht ändern!
	filecolor=black,					% bitte nicht ändern!
	urlcolor=black,						% bitte nicht ändern!
	citecolor=black,					% bitte nicht ändern!
	pdftex=true,              % bitte nicht ändern!
	plainpages=false,         % bitte nicht ändern!
	hypertexnames=false,      % bitte nicht ändern!
	pdfpagelabels=true,       % bitte nicht ändern!
	hyperindex=true]{hyperref}% bitte nicht ändern!
\else
  % hier kann mann eventuelle Befehle umdefinieren
  % die nur für pdfLaTeX vorgesehen sind
  % und das richtige Kompilieren durch den normalen LaTeX verhindern
	\newcommand{\texorpdfstring}[2]{#2}
\fi

%%%%%%%%%%%%%%%%%%%%%%%
% Hyphenation options %
%%%%%%%%%%%%%%%%%%%%%%%

%For manual hypenation settings: \hyphenation{Hy-phenation} etc.

%%%%%%%%%%%%%%%%%%%
% Listing options %
%%%%%%%%%%%%%%%%%%%

%Define pseudo-code
\lstdefinelanguage{pseudo-code}{texcl,mathescape,keywords={function,begin,end,return,until,repeat,while,do,case,for,each,if,then,else,or,endcase,loop,od,continue_while,continue_loop,when,wait,from,in,send,to,not,union,downto,double,boolean,this,new,class}, comment=[l]{//},breaklines=true,breakatwhitespace=true}

%Define basic options
\lstset{frameround=tttt,basicstyle=\small}




\setcounter{secnumdepth}{4}
\setcounter{tocdepth}{4}
\renewcommand{\thechapter}{\Roman{chapter}}
\renewcommand{\thesection}{\arabic{section}}
%\renewcommand{\thefigure}{\arabic{figure}}




\begin{document}
\setstretch{1.2} % Zeileabstand: NICHT über 1.5 setzen!!

%Front page and abstract
\thispagestyle{empty}
\vspace{2cm}
\begin{center}
\includegraphics[scale=1]{images/uni-logo} \hspace{3cm} 
\includegraphics[scale=1]{images/uni-logo} % Logo des Betriebs
\\\vspace{2cm}

\Huge{\dctitle}

\vspace{1cm}
\Large{Bachlorarbeit}

\vspace{5mm}
\large{\dcfaculty\\ \dcuniversity}
\vspace{3.5	cm}
\end{center}
{\large
\begin{flushleft}
%\linespread{3}
\textbf{Vorgelegt von:} \dcauthorsurname, \dcauthorname \\
\textbf{Geboren am:} \dcauthoradd\\
\textbf{Matr.-Nr.:} \dcauthorno \\
\textbf{Gutachter:} \\
\begin{tightlist}
	\item \dcapprovala
	\item \dcapprovalb
\end{tightlist}

\textbf{Betrieblicher Betreuer:} \\
\begin{tightlist}
	\item \dcada
\end{tightlist}
\vspace{5mm}
\textbf{Abgabedatum:} \dcdatesubmitted\\
\end{flushleft}
}
\vspace{1cm}

\newpage



\tableofcontents

%Index
\clearpage
\label{sec:index}
\printindex

%Content



%\input{introduction}
%\input{background}
\chapter{Einleitung}


\begin{flushleft}\end{flushleft}

\chapter{Analyse}

\section{Einordnung des Kreditorenworkflows}
Um die Intention der Neukonzeption verstehen und sie angemessen durchführen zu können, ist es zunächst wichtig, den Workflow der Kreditorenbuchhaltung im Unternehmen einordnen zu können. 
Dieser Umstand ist analog zu den Anforderungen der ISO 9000 Reihe zu sehen, das zur Qualitätssicherung in Prozessen ein Verständnis von deren Abhängigkeiten und Wechselwirkungen vorgibt. 
Auf diese Weise soll ersichtlich werden, welche Anforderungen gegenüber dem Workflow bestehen.

\subsection{Prozessportfolio}
Die ISO 9000 fordert für
Die Einordnung ein ein Prozessportfolio gemäß Six Sigma.



\subsection{Prozesslandkarte}
$\rightarrow$ Nicht klassischer Unterstützungsprozess\\
$\rightarrow$ Eher Kernprozess



\subsection{Auswertung der Einordnung}





\section{Ist-Aufnahme}
Der Prozess wird in zwei sequentiellen Stufen betrachtet. 
Im ersten Teil geht die Rechnung  ein und wird freigegeben.
Im zweiten Teil wird die Rechnung gebucht und der Zahllauf vorbereitet.
\subsection{Kreditorenworkflow Teil 1: Rechnungsfreigabe}
Der erste Teil beginnt streng genommen mit der Kostenverursachung, auf die eine Kreditorenrechnung folgt. 
Da die Kostenverursachung allerdings für den Workflow nicht weiter relevant ist, wird der Rechnungseingang als Ausgangspunkt des Prozesses und Auslöser einer Prozessinstanz betrachtet. 
\subsubsection{Prozessablauf}
Die eingehende Kreditorenrechnung wird innerhalb des Unternehmens an einen bestimmten Mitarbeiter der Finanzbuchhaltung gegeben.
Für die nachfolgenden Schritte gibt es zwei mögliche Szenarien.

\paragraph{Fall 1: Kostenverursacher in Niederlassung}
~\\ 
In diesem Szenario befindet sich der Kostenverursacher nicht in der Unternehmenszentrale, also in räumlicher Distanz zur Finanzbuchhaltung.\\
Der Mitarbeiter scannt die Rechnung ein, um das Original digital zu konservieren, was den ersten Schritt der Versionierung darstellt. $\rightarrow$
Danach muss die Rechnung zur sachlichen Freigabe an den Kostenverursacher gegeben werden. 
Wird die Rechnung sachlich nicht freigegegeben, endet die Prozessinstanz.
Nachfolgend müssen vom Mitarbeiter der Finanzbuchhaltung die Gründe dafür eruiert und ggf. eine korrekte Rechnung nachgefordert werden.$\rightarrow$
Wird die Rechnung sachlich freigegeben, druckt der freigebende Mitarbeiter dazu die Rechnung aus, unterschreibt und scannt wieder.
Das Resultat sendet er dem Mitarbeiter der Finanzbuchhaltung.
Dieser druckt wiederum das Resultat aus.
Danach muss die fachliche Freigabe erfolgen.
Dazu verbringt der Mitarbeiter die Original-Rechnung und den Ausdruck mit der sachlichen Freizeichnung an den verantwortlichen Bereichsleiter.
Dieser lehnt ggf. die Rechnung ab, dazu wird wie im Fall der sachlichen Freigabe verfahren.
Im Regelfall gibt er mit seiner Unterschrift die Rechnung fachlich frei.
Ggf. erfolgt ein Vermerk, dass die Rechnung an einen Kunden weiterzuberechnen ist. 
Nachdem die fachliche Freigabe erfolgt ist, wird die finale Version eingescannt.\\
Damit liegen alle Freigaben vor und der Mitarbeiter ist zur Buchung berechtigt.

\paragraph{Fall 2: Kostenverursacher in Unternehmenszentrale}
~\\
In diesem Szenario befindet sich der Kostenverursacher in der Unternehmenszentrale, also in für die Mitarbeiter der Finanzbuchhaltung erreichbarer Nähe.
Zunächst muss ebenso die sachliche Freigabe erfolgen.
Dazu verbringt der Mitarbeiter der Finanzbuchhaltung die Original-Rechnung an den Kostenverursacher.
Für eine nicht erfolgte Freigabe wird wie in Fall 1 verfahren.
Bei erfolgter sachlicher Freigabe hat ebenso als nächtes die fachliche Freigabe zu erfolgen.
Wie in Fall 1 wird dazu die Original-Rechnung, diesmal allerdings mit sachlicher Freigabe an den Bereichsverantwortlichen verbracht.
Auch an dieser Stelle wird wie in Fall 1 für Freigabe oder Nicht-Freigabe verfahren.
Ggf. erfolgt auch hier ein Vermerk, der die Weiterberechnung anweist.
Im Regelfall gibt er mit seiner Unterschrift die Rechnung fachlich frei.
Abschließend wird die vollständig freigezeichnete Version eingescannt.
Damit liegen alle Freigaben vor und der Mitarbeiter ist zur Buchung berechtigt.

\subsubsection{Prozessmodellierung}
Der Prozess lässt sich soweit modellieren. 
Zu beachten sind dabei zwei Prozessschnittstellen. 
Eine davon ist die zur Debitorenbuchhaltung, die bei einem Weiterberechnungsvermerk greift. 
Diese wird modelliert, ist allerdings inhaltlich in diesem Rahmen nicht relevant.
Die zweite Schnittstelle ist die Verbindung zum zweiten Teil des Workflows.\\
Auf die Unterscheidung zwischen der internen Verbringung der Rechnung und der Versendung per Mail wird in dieser Modellierung verzichtet, da sie im Ablauf keine inhaltlich entscheidenden Unterschiede beinhaltet, sondern nur das Papieraufkommen erhöht.
\\
DAS MUSS ABER IN DEN ANHANG


\subsection{Kreditorenworkflow Teil 2: Rechnungsbuchung}
Im zweiten Teil werden zur Buchung bereits Bestands-EDV-Systeme benutzt.
\subsubsection{Prozessablauf}
Zur Buchung werden alle eingescannten Versionen inklusive des Originals verwendet.
Zur Buchung wird der Beleg an Datev Rechnungswesen Pro übergeben.
Das Feature der automatischen Zeichenerkennung erkennt dabei insbesondere die Rechnungsnummer, das Datum, den Saldo und die Steuer.
Die Kostenstelle, sowie der Kostenträger müssen durch den Mitarbeiter angegeben werden.
In Datev wird die Rechnung dem Kreditor und unter diesem wiederum der Rechnungsnummer zugeordnet.
Damit ist die Rechnung in Datev verbucht.\\
Die Rechnung muss anschließend noch gezahlt werden.
Dazu wird ein Zahllauf in Datev angelegt.
In diesem erscheinen die fälligen Belege. 
Die Fälligkeit ergibt sich hierbei aus Rechnungsdatum und Zahlungsziel, welches aus den Kreditorenstammdaten entnommen wird.
Die fälligen Zahlungen werden nun zunächst durch den Mitarbeiter selbst geprüft.
Anschließend wird der Zahllauf dem Vorstandsvorsitzenden in Schriftform zur Freigabe vorgelegt.
Ist die Freigabe erfolgt, werden die in einer Datei an das eBanking-System gegeben.
Ein zweiter Mitarbeiter prüft die Freigabe durch den Vorstandsvorsitzenden und gibt die Zahlungen im eBanking-System frei.
Das eBanking-System übermittelt daraufhin die Zahlungen an die Bank, wo die Zahlung durchgeführt wird.
\subsubsection{Prozessmodellierung}

Auch dieser Teilprozess lässt sich mit dem gegebenen Ablauf modellieren. 
Die Durchführung der Zahlung kann als transparent betrachtet werden, da sie vollautomatisch durchgeführt wird.





\section{Problemanalyse}
Die nun folgende Problemanalyse ist Ausgangspunkt für die Neukonzeption.

\subsection{Methodik der Problemanalyse}
Zur Identifikation von problematischen Aspekten wird zunächst die argumentative Grundlage geschaffen, diese zweifelsfrei ermitteln zu können.
 dienen einerseits die Prozessbeschreibungen bzw. Prozessmodelle und andererseits die Interviews, die 
\subsubsection{Stakeholderanalyse}
Die Stakeholder des Prozesses ergeben sich aus den Akteuren, die in der Prozessbeschreibung bzw. der Prozessmodellierung als Beteiligte auftreten.\\
Das sind:


\begin{enumerate}
\item{Mitarbeiter der Finanzbuchhaltung}
\\ Speziell der Mitarbeiter der Finanzbuchhaltung, der für die Kreditorenbuchhaltung zuständig ist, ist in diesem Rahmen aktiv. 
\\ Ein weiterer Akteur aus diesem Bereich ist der (faktisch wechselnde) Mitarbeiter, der den Zahllauf nach dem Vieraugen-Prinzip freigibt. 
Dieser ist jedoch nur marginal beteiligt.
\\ Ein Stakholder, der nicht maßgeblich in diesem Prozess agiert, ist der Leiter der Finanzbuchhaltung, der die Prozesse der Abteilung plant, deren Implementation steuert und den Ablauf kontrolliert.

\item{Kostenstellenverantwortliche}
\\ Die Kostenverursacher sind Mitarbeiter, die zur Kostenverursachung befugt sind. 
Faktisch sind dies Kostenstellenverantwortliche, die regelmäßige oder abwicklungsbezogene Kosten verursachen und damit Kreditorenrechnungen auslösen. 
Durch sie erfolgt die sachliche Freigabe.

\item{EDV-Abteilung}
\\ Da der Prozess teilweise technologiegestützt abläuft, ist die Abteilung der Informationstechnologie, wenn auch nicht aktiv an Prozessinstanzen beteiligt, als Stakeholder auszumachen. 

\item{Bereichsverantwortliche}
\\ Die Mitarbeiter, die die fachliche Freigabe erteilen dürfen, sind Bereichsleiter. 
Konkret lässt sich dies zum Zeitpunkt der Erstellung dieser Arbeit auf drei Personen konkretisieren. 
Von diesen befinden sich zwei im Vorstand, der dritte ist der Bereichsleiter Informationstechnologie.\\

\end{enumerate}

Aus dieser Liste lassen sich konkrete Personen ableiten, deren Erfahrung mit dem Prozess zur Problemidentifikation dienen kann. 
Mit diesen werden Interviews geführt, um eine empirische Grundlage zu schaffen.

Es ergeben sich:

\begin{enumerate}
\item{Mitarbeiter der (Kreditoren-)Buchhaltung}
\item{Leiter der Finanzbuchhaltung}
\item{Ein Bereichsverantwortlicher}
\item{Leiter der Informationstechnologie}
\end{enumerate}


\subsubsection{Beschreibung und Modellierung der Teilprozesse}

Auch die erarbeiteten Abläufe der Teilprozesse dienen als Grundlage der Analyse.


\subsection{Fehler}
Eine von zwei Möglichkeiten, die im Zusammenhang mit dem Prozess problematisch sind, sind fehlerhafte Prozessinstanzen. Als fehlerhaft wird eine Prozessinstanz betrachtet, wenn diese nicht wie beschrieben abläuft, also der Modellierung nicht folgt.
\subsubsection{Fehlerarten}

\begin{enumerate}
\item{Rechnung nicht zugegen}
\\ Die Rechnung befindet sich zu einem bestimmten Zeitpunkt nicht am vorgesehenen Ort. 
Sie ist entweder nicht in die Finanzbuchhaltung gegeben worden oder auf sonstige Art auf dem Weg des Freigabeprozesses verloren gegangen.
\item{Formal nicht korrekte Rechnung}
\\ Die Rechnung weist nicht alle Daten auf, die zur sachlichen und fachlichen Freigabe ggf. nötig wären, sodass Rücksprache mit dem Kreditor notwendig ist.
\end{enumerate}


\subsubsection{Fehlerquellen}

Die Fehlerquellen sind aus den Interviews ersichtlich. 
Die bereitgestellten Systeme werden als adäquat und funktional betrachtet.
Als problematisch wird nur der Teil empfunden, in dem keine IT-System-seitige Prozessunterstützung stattfindet.
Die Gründe für fehlerhafte Prozessinstanzen sind demnach auf menschliches Verhalten zurückzuführen, welches in diesem Zusammenhang durch Nachlässigkeit die beschrieben Fehler auslöst.
\\
Bezüglich fehlerhafter Rechnungen ist die Ursache meistens auf Nachlässigkeit seitens des Kreditors zurückzuführen.
Da diese von internen Prozessen losgelöst ist, steht diese Fehlerquelle nicht im Fokus dieser Arbeit. 

\subsection{Performance}
Die zweite Möglichkeit für eine als problematisch wahrgenommene Prozessinstanz ist deren mangelhafte Performance.
Diesbezüglich gibt es mehrere Anhaltspunkte, die von Stakeholdern angesprochen werden.
Zum einen gestalten sich Prozessdurchläufe langwierig.
Unabhängig davon, ob die Freizeichnung mit Akteuren durchgeführt wird, die sich inner- oder außerhalb der Unternehmenszentrale befinden und dementsprechend eine räumliche Entfernung zwischen der Finanzbuchhaltung und den Freizeichnern ergibt vorliegt oder nicht, ist der Freigabeprozess mit hohem manuellen Kommunikationsaufwand verbunden, der über E-Mail oder direkt, allerdings mit Laufwegen durch das Unternehmen bearbeitet wird.
\\
Aus diesem Umstand ergibt sich auch die mangelhafte Kontrollmöglichkeit, sowohl für den betroffenen Mitarbeiter der Finanzbuchhaltung, sowie die mit der Freizeichnung beauftragten Akteure, da sich Rechnungen in der zu einem bestimmten Zeitpunkt aktuellen Version faktisch nur im Zugriff eines einzelnen Mitarbeiters befinden. 
Besonders deutlich wird diese Situation bei Rechnungen, die innerhalb der Unternehmenszentrale durchgereicht werden.
An jeder Station (also Rechnungsannahme, sachliche Freigabe, fachliche Freigabe und wieder zurück in der Finanzbuchhaltung) besteht zusätzlich die Möglichkeit, dass ein Beleg abhanden kommt und eine neue Prozessinstanz erst durch Nachsendung einer Mahnung oder Zahlungserinnerung ausgelöst wird.

\subsubsection{Key Performance Indicators}
Diese Kritikpunkte lassen sich in objektiven Kennzahlen ausdrücken, an denen der Prozess im Ist-Zustand gemessen werden kann. 
Diese sollen auch die Neukonzeption des Prozesses begleiten, bzw. diese an der Verbesserung resultierender Werte ausgerichtet werden.
Auf Basis der Erläuterungen in 3.2 und in diesem Kapitel ergeben sich zwei Kennzahlen:

\begin{enumerate}
\item{Prozesslaufzeit}
\\ Die objektiv messbare Zeit von Rechnungseingang/-annahme bis Bezahlung, speziell jedoch in dem als performanceschwach wahrgenommenen Teil von Rechnungseingang/-annahme bis Buchung.
\item{Prozessfehlerrate}
\\ Relativer Anteil fehlerhafter Prozessinstanzen.
\end{enumerate}




\begin{flushleft}\end{flushleft}

\chapter{Neukonzeption}

\section{ADM als Methode}
Die Architecture Development Method (ADM) ist der Kernbestandteil von TOGAF. 
ADM ist ein iteratives Modell zur Konzeption und Verwaltung von Geschäftsarchitekturen, das aus mehreren Phasen besteht.
Zentrales Element dieses Modells ist das Anforderungsmanagement.
Dieses sagt aus, dass sich die Schritte entlang des Zyklus immer wieder an den Anforderungen und geschäftlichen Zielen ausgerichtet werden, welche sich im Laufe eines Projekts z.B. durch geschaffene Möglichkeiten verändern können.\\
-> Bild ADM
\\
Dabei gilt, dass jede Phase einen Output liefert, der als Input für die nächste Phase benutzt wird.\\
Aufgrund des Umfanges werden in dieser Arbeit die Phasen bis C betrachtet. Die konzeptionelle Arbeit wird bis Phase D geleistet, diese wird allerdings ausgeklammert.
\begin{enumerate}
\item{Preliminary}
\\ In der vorgelagerten Phase findet die Eingrenzung des Umfangs auf Basis der Anforderungen statt. 
Ein weiterer Bestandteil ist das "Tailoring" statt, in dem das methodische Vorgehen geplant und ggf. gegenüber dem Framework (TOGAF) angepasst wird.\\
Output ist das Dokument "Request For Architecture Work", das Ausgangspunkt des Requirement Management ist.
\item{Phase A: Architecture Vision}
\\ In dieser Phase wird richtungsweisend definiert, auf welche Art die zu konzeptionierende Infrastruktur den Anforderungen und Zielen Rechnung tragen soll. Dabei wird auch auf Weiter- oder Wiederverwendungsmöglichkeiten von Bestandsressourcen geschaut.\\
Dies wird in das Dokument "Statement Of Architecture Work" als Output überführt.
\item{Phase B: Business Architecture}
\\ Phase B definiert die Geschäftsarchitektur hinsichtlich Aufbau- und Ablauforganisation. Das bedeutet, dass hierbei die Anforderungen in konkrete Funktionen und Abläufe überführt werden, was den Output dieser Phase darstellt.
\item{Phase C: Information Systems Architectures}
\\ Die letzte betrachtete Phase definiert die technische Architektur hinsichtlich Datenhaltung und -verarbeitung (Application Architecture). \\
Ein Datenhaltungsmodell und die erforderliche Anwendungs- und Schnittstellenlandschaft werden festgelegt.


\end{enumerate}




\section{Preliminary}
\subsection{Anforderungsanalyse}
\subsubsection{Interne Anforderungen}
\subsubsection{Gesetzliche Anforderungen}

Wird in TOGAF als Preliminary bezeichnet.

\section{Phase A: Architecture Vision}
Phase A Architecture Vision
\section{Phase B: Business Architecutre}
Phase B Business Architecture

\section{Phase C: Information Systems Architectures}
Phase C Information System Architectures

\subsection{Data Architecture}

\subsection{Application Architecture}

\section{Technologieinfrastruktur}
\begin{flushleft}\end{flushleft}

\chapter{Zusammenfassung}




\chapter{Anhang}
%\input{results}
%\input{conclusions}

%Bibliography
\bibliographystyle{apalike-german} % Kann auch gerne ausgetauscht werden; plain spart z.B. Platz
\bibliography{literature}

%Glossary
%See e.g. http://www.bastian.name/blogger/2006/02/abkrzungsverzeichnis-und-glossar.html for further options and how to integrate this with texniccenter
\clearpage
%
\glossary{name={Diplomarbeit},description={Eine Diplomarbeit ist eine wissenschaftliche Arbeit, die zur Erlangung des Abschlusses Diplom an einer Universit�t, Fachhochschule oder Berufsakademie geschrieben werden muss.(Wikipedia)}}

%\printglossary
\automark[section]{chapter}

%List of figures
\clearpage
\listoffigures


%Appendices
%\input{appendices}

%Adding formal stuff...


\end{document}
